\chapter{API Parasitism: Exploiting Web2 Frailties}

\section{Introduction}
As Web3 infrastructure matures, its interaction with legacy Web2 systems becomes increasingly necessary. However, Web2 APIs are notoriously fragile, prone to rate-limiting, downtime, and arbitrary deprecation. API Parasitism is the architectural methodology of reliably extracting data from these hostile environments.

\section{Decentralized Oracle Architecture}
A robust Decentralized Oracle Network (DON) must abstract the fragility of the underlying Web2 APIs. The architecture involves an array of independent nodes, each querying multiple data sources. The results are aggregated using Byzantine Fault Tolerant consensus algorithms to derive a single, verifiable truth. This redundancy ensures that the failure of any single API endpoint does not compromise the integrity of the smart contract relying on the data.

\section{JSON Parsing Fallbacks and Data Extraction}
Web2 APIs frequently alter their JSON response schemas without notice. A resilient oracle must implement dynamic JSON parsing fallbacks. This involves utilizing fuzzy matching and heuristic analysis to extract critical data fields even when the expected JSON path is broken. By mapping semantic relationships within the data payload, the parser can adapt to structural mutations in real-time.

\section{Specific Web2 API Frailties}
The inherent weaknesses of Web2 APIs are well-documented. For instance, the infamous Twitter API rate limit changes of 2023 demonstrated the danger of centralized data monopolies. Furthermore, RESTful architectures often lack built-in cryptographic signatures, making them susceptible to man-in-the-middle attacks. By leveraging API Parasitism, decentralized applications can treat these APIs not as trusted partners, but as unreliable sensors in a broader, fault-tolerant network.
